% Content of the blueprint, shared by the web and the print version.
% Each chapter follows one section of the paper; the chapter files are in chapters/.
% The paper is cited by its own statement numbers and LaTeX labels.

\section*{The paper}

This blueprint follows the paper \emph{Lions' maximal regularity problem for divergence-form
differential operators: failure at the $\frac12$-H\"older endpoint} by Lukas Niebel
(arXiv:2608.11194v2) section by section: Section~1 (introduction, Theorem~1.1 and
Corollary~1.2), Section~2 (spatial oscillatory modes), Section~3 (time blocks, frequencies and
the coefficient), Section~4 (the variational solution and its time derivative), Section~5
(cancellation and the forcing), Section~6 (full-space and bounded-domain consequences) and
Section~7 (comparison with endpoint hypotheses in the literature).  A last chapter records the
part of Lions' theory that the formalisation proves, which the paper cites.  The paper is cited
by its own statement numbers and LaTeX labels --- ``Theorem~1.1 of the paper'',
``\texttt{eq:reciprocal-sums}'' --- and a \texttt{\textbackslash ref} inside the blueprint
prints the \emph{blueprint's} number, which is in general not the paper's.

\section*{Reading guide}

\textbf{What the green marks mean.}  \texttt{\textbackslash leanok} on a \emph{statement}
means that the Lean declarations named in that node's \texttt{\textbackslash lean} list exist in
this repository, elaborate, and together say what the statement says; for a
\texttt{definition} node it means that the Lean definition exists and formalises the object
described.  \texttt{\textbackslash leanok} inside a \emph{proof} means that the Lean proofs of
those declarations are complete --- no \texttt{sorry}, and no hypothesis that the paper does not
also have or discharge.  Every statement of the blueprint carries a \texttt{\textbackslash lean}
list and \texttt{\textbackslash leanok}, and so does every proof.  This includes every numbered
statement of the paper except Remark~3.1 (\emph{Choice of the frequencies}), which makes no
claim used later and has no Lean counterpart.  Theorem~\ref{thm:main} is
\texttt{LionsC12.lions\_counterexample\_interval} and Corollary~\ref{cor:bounded-domains} is
\texttt{LionsC12.lions\_counterexample\_domain}; both have the statements of
\texttt{Challenge.lean}.

\textbf{Where the differences are recorded.}  Every node whose Lean form differs from the
paper's --- an explicit constant in place of ``$C$'', a pointwise estimate in place of an
$\Leb^\infty$ one, a different route --- is followed by a paragraph \textbf{Lean form}, outside
the theorem environment, that says exactly how.  Those paragraphs, not the statements, are where
the differences are recorded; a green statement with no such paragraph is formalised as written.
\texttt{DEVIATIONS.md} collects them.

\chapter{Introduction: the theorems and their Lean statements}
The paper works with the complex Gelfand triple
\[
  V=\Sob_0^1((0,\pi);\CC)\hookrightarrow H=\Leb^2((0,\pi);\CC)\hookrightarrow V'=\Sob^{-1}((0,\pi);\CC),
\]
$V'$ the anti-dual of $V$, and with the form operators $\cA(t):V\to V'$ of a bounded real
coefficient $\fra$.  This chapter introduces the objects of the statement layer
(\texttt{LionsC12/Paper/Statement.lean}, repeated character for character in
\texttt{Challenge.lean}) and states the two main results.  Everything is built from Mathlib: no
Sobolev space, anti-dual or Bochner space is assumed.

\section{The function spaces}

\begin{definition}[The spaces $H$, $V$ and $V'$ on an open subset of the line]
  \label{def:spaces}
  \lean{LionsC12Challenge.L2, LionsC12Challenge.testGraph, LionsC12Challenge.H10Graph,
    LionsC12Challenge.H10, LionsC12Challenge.val, LionsC12Challenge.dx,
    LionsC12Challenge.Hm1, LionsC12Challenge.ofL2, LionsC12Challenge.ofH10}
  \leanok
  For an open $U\subseteq\R$, $\Leb^2(U;\CC)$ is Mathlib's $\Leb^2$ space of
  $\mathrm{volume}|_U$.  $\Sob_0^1(U;\CC)$ is the closure, in the Hilbert product
  $\Leb^2(U)\times\Leb^2(U)$ with norm $(\norm{v}^2+\norm{v'}^2)^{1/2}$, of the pairs
  $(\varphi,\varphi')$ with $\varphi\in\C_c^1(\R;\CC)$, $\supp\varphi\subseteq U$ --- the graph
  definition of $\Sob_0^1$ (Brezis).  The two projections are $v\mapsto v$ (\texttt{val}) and
  $v\mapsto\partial_xv$ (\texttt{dx}).  $V'=\Sob^{-1}(U;\CC)$ is the space of continuous
  conjugate-linear functionals on $V$, with its operator norm; $H\hookrightarrow V'$ is
  $h\mapsto(w\mapsto\langle h,w\rangle_H)$ and $V\hookrightarrow V'$ is the composite.
\end{definition}

\textbf{Lean form.}  \texttt{H10 U} is a type synonym of the closed submodule
\texttt{H10Graph U}, with its Hilbert-space structure declared explicitly; this avoids a
type-class search that does not terminate for the anti-dual of a submodule.  The inner product
of $V$ is $\langle v,w\rangle_H+\langle\partial_xv,\partial_xw\rangle_H$.

\begin{definition}[The forms and the variational solutions]
  \label{def:lions-solution}
  \lean{LionsC12Challenge.form, LionsC12Challenge.HasWeakDeriv, LionsC12Challenge.IsLionsSolution}
  \leanok
  \uses{def:spaces}
  For a coefficient $\fra:\R\times\R\to\R$,
  \[
    \fra(t)[v,w]=\int_U\fra(t,x)\,\partial_xv(x)\,\overline{\partial_xw(x)}\dd x
    =\langle\cA(t)v,w\rangle_{V',V}.
  \]
  A function $g:(0,1)\to E$ is the distributional derivative of $u:(0,1)\to E$ if
  $\int_0^1\varphi'u\dd t=-\int_0^1\varphi g\dd t$ for every real $\varphi\in\C_c^\infty(0,1)$.
  A \emph{Lions variational solution} with zero initial value of $u'+\cA(t)u=f$ is a function
  $w\in\Leb^2(0,1;V)$ whose $V'$-valued distributional derivative $w'$ lies in
  $\Leb^2(0,1;V')$, such that $w'(t)+\cA(t)w(t)=f(t)$ in $V'$ for almost every $t$, and whose
  continuous representative in $\C([0,1];V')$ vanishes at $t=0$.
\end{definition}

\textbf{Lean form.}  $\Leb^2(0,1;X)$ is \texttt{MemLp $\cdot$ 2 (volume.restrict (Ioo 0 1))},
and the Bochner integrals are Mathlib's; the equation in $V'$ is tested against every $w\in V$.
The initial condition is imposed on a representative \texttt{w₀} continuous on $[0,1]$ and equal
to $w$ almost everywhere, which is how ``$w(0)=0$'' is read in $W(0,1)\subseteq\C([0,1];H)$.

\begin{definition}[The spaces on a domain of $\R^d$]
  \label{def:spaces-domain}
  \lean{LionsC12Challenge.L2Ω, LionsC12Challenge.L2Ωvec, LionsC12Challenge.grad,
    LionsC12Challenge.testGraphΩ, LionsC12Challenge.H10ΩGraph, LionsC12Challenge.H10Ω,
    LionsC12Challenge.valΩ, LionsC12Challenge.gradΩ, LionsC12Challenge.Hm1Ω,
    LionsC12Challenge.ofL2Ω, LionsC12Challenge.ofH10Ω, LionsC12Challenge.isoCoeff,
    LionsC12Challenge.formΩ, LionsC12Challenge.IsLionsSolutionΩ}
  \leanok
  \uses{def:spaces, def:lions-solution}
  For an open $\Omega\subseteq\R^d$ the same constructions with
  $(\varphi,\nabla\varphi)\in\Leb^2(\Omega;\CC)\times\Leb^2(\Omega;\CC^d)$:
  $H_\Omega=\Leb^2(\Omega;\CC)$, $V_\Omega=\Sob_0^1(\Omega;\CC)$, $V_\Omega'=\Sob^{-1}(\Omega;\CC)$, the
  isotropic matrix $\frB(t,x)=\frb(t,x)\Id_d$ as a continuous linear map of $\R^d$, the form
  $\frb(t)[v,w]=\int_\Omega\frb(t,x)\,\nabla v\cdot\overline{\nabla w}\dd x$, and the Lions
  solutions of $U'-\Div(\frB\nabla U)=F$.
\end{definition}

\begin{proposition}[$\Sob_0^1$ through weak derivatives]
  \label{prop:sobolev}
  \lean{LionsC12.HasWeakDerivL2, LionsC12.W12Graph, LionsC12.isClosed_W12Graph,
    LionsC12.hasWeakDerivL2_unique, LionsC12.H10Graph_le_W12Graph,
    LionsC12.H10Graph_eq_closure_smoothGraph, LionsC12.range_H10toW12,
    LionsC12.HasWeakGradL2, LionsC12.W12ΩGraph, LionsC12.isClosed_W12ΩGraph,
    LionsC12.hasWeakGradL2_unique, LionsC12.H10ΩGraph_le_W12ΩGraph,
    LionsC12.H10ΩGraph_eq_closure_smoothGraphΩ, LionsC12.range_H10ΩtoW12Ω}
  \leanok
  \uses{def:spaces, def:spaces-domain}
  Let $U\subseteq\R$ be open, and let $\W^{1,2}(U)$ be the space of the pairs
  $(v,g)\in\Leb^2(U)\times\Leb^2(U)$ with $\int_Uv\varphi'\dd x=-\int_Ug\varphi\dd x$ for every
  real $\varphi\in\C_c^\infty(U)$, a closed subspace on which $g$ is determined by $v$.  Then
  $\Sob_0^1(U)\subseteq\W^{1,2}(U)$, and $\Sob_0^1(U)$ is the closure of $\C_c^\infty(U)$ in
  $\W^{1,2}(U)$.  The same holds on an open $\Omega\subseteq\R^d$, with weak gradients.
\end{proposition}

\textbf{Lean form.}  Not in the paper, which uses the textbook definitions; it identifies the
graph definition of Definition~\ref{def:spaces} with them.  $\W^{1,2}(U)$ is the submodule
\texttt{W12Graph U} of $\Leb^2(U)\times\Leb^2(U)$, with the norm $(\norm{v}^2+\norm{g}^2)^{1/2}$;
the inclusion $V\hookrightarrow\W^{1,2}(U)$ is an isometry whose range is the closure of the
graphs of the $\C_c^\infty(U)$ functions (\texttt{range\_H10toW12}).

\begin{proof}
  \leanok
  The inclusion is integration by parts against the test graphs.  A $\C_c^1$ function supported in
  $U$ is, with its derivative, the uniform limit on a fixed compact subset of $U$ of its
  convolutions with smooth bumps of small radius, which are $\C_c^\infty$ functions supported in
  $U$; uniform convergence on a set of finite measure gives convergence in $\Leb^2$.
\end{proof}

\section{The main results}

\begin{theorem}[Counterexample on an interval]
  \label{thm:main}
  \lean{LionsC12.lions_counterexample_interval}
  \leanok
  \uses{def:lions-solution}
  Let $\delta>0$.  There exists a bounded, uniformly elliptic, real-valued coefficient $\fra$ with
  \[
    \fra\in\C^{0,1/2}([0,1];\Leb^\infty(0,\pi)),\qquad
    \norm{\fra-1}_{\C^{0,1/2}([0,1];\Leb^\infty(0,\pi))}<\delta,
  \]
  but $\fra\notin\C^{0,\alpha}([0,1];\Leb^\infty(0,\pi))$ for every $\frac12<\alpha<1$.  There also
  exist a real-valued forcing $f\in\C([0,1];\Leb^2(0,\pi))\cap\Leb^2(0,1;\Leb^2(0,\pi))$ and a
  real-valued $u\in\C([0,1];V)\cap\C^1([0,1];V')\cap\W^{1,1}(0,1;H)$ such that $u(0)=0$ and $u$
  is the unique Lions variational solution of $\partial_tu(t)+\cA(t)u(t)=f(t)$ in $V'$ for every
  $t\in[0,1]$.  Moreover $u(t)\in D(A_H(t))$ for every $t\in[0,1]$, and
  \[
    A_H(\cdot)u(\cdot)\in\Leb^1(0,1;H),\qquad A_H(\cdot)u(\cdot)\notin\Leb^2(0,1;H),\qquad
    \partial_tu\notin\Leb^2(0,1;H).
  \]
\end{theorem}

\textbf{Reading of the clauses in Lean.}  The statement is that of
\texttt{LionsC12Challenge.lions\_counterexample\_interval} in \texttt{Challenge.lean}, proved in
\texttt{Solution.lean} by \texttt{LionsC12.lions\_counterexample\_interval}.  The $\C^{0,1/2}$ norm is
$\sup_t\norm{g(t)}+\sup_{s\neq t}\norm{g(t)-g(s)}/\abs{t-s}^{1/2}$, so the smallness is the
existence of bounds $A,B\geq0$ for the two suprema with $A+B<\delta$; the $\Leb^\infty$ norms
are \texttt{eLpNorm $\cdot$ ∞} on $(0,\pi)$.  $u\in\C^1([0,1];V')$ is the existence of a
continuous $u'$ with \texttt{HasDerivWithinAt} at every $t\in[0,1]$, relative to $[0,1]$.
$u\in\W^{1,1}(0,1;H)$ is Bochner integrability of $u$ in $H$ together with an integrable
$H$-valued distributional derivative.  $D(A_H(t))$ is the set of $v\in V$ with $\cA(t)v\in H$.
``$A_H(\cdot)u(\cdot)\notin\Leb^2$'' is asserted of \emph{every} $g$ with
$\cA(t)u(t)=g(t)$ for all $t$, and ``$\partial_tu\notin\Leb^2(0,1;H)$'' says that no
$g\in\Leb^2(0,1;H)$ represents $u'$ almost everywhere.  Real-valuedness of an element of
$\Leb^2$ is the vanishing of its imaginary part almost everywhere.  Uniqueness is proved for
every Lions solution in the sense of Definition~\ref{def:lions-solution}.

\begin{proof}
  \leanok
  \uses{lem:oscillatory-modes, prop:coefficient, prop:holder-sharp, lem:solution-regularity,
    prop:forcing, lem:lions-uniqueness, lem:domain-membership}
  Take $\eps=\min\{1,\delta/(2(1+2C_\beta))\}$ and $\fra=1+\eps\beta$.
  Proposition~\ref{prop:coefficient} gives the bounds $A=\eps/480$ and
  $B=2\eps C_\beta$, Proposition~\ref{prop:holder-sharp} the failure of $\alpha$-H\"older
  continuity, Lemma~\ref{lem:solution-regularity} the regularity of $u$ and
  $\partial_tu\notin\Leb^2$, Proposition~\ref{prop:forcing} the forcing and the equation for
  every $t$, Lemma~\ref{lem:domain-membership} the clauses on $A_H(\cdot)u(\cdot)$, and
  Lemma~\ref{lem:lions-uniqueness} uniqueness.
\end{proof}

\begin{corollary}[Counterexamples on arbitrary bounded domains]
  \label{cor:bounded-domains}
  \lean{LionsC12.lions_counterexample_domain, LionsC12.lions_counterexample_open}
  \leanok
  \uses{def:spaces-domain}
  Let $d\geq1$, let $\Omega\subset\R^d$ be a bounded domain and $\delta>0$.  There is a bounded
  real scalar coefficient $\frb$ such that $\frB=\frb\Id_d$ is real symmetric, uniformly
  elliptic, $\frB\in\C^{0,1/2}([0,1];\Leb^\infty(\Omega;\R^{d\times d}))$ with
  $\norm{\frB-\Id_d}_{\C^{0,1/2}}<\delta$, and $\frB\notin\C^{0,\alpha}$ for every
  $\frac12<\alpha<1$.  There are a real forcing $F\in\C([0,1];H_\Omega)\cap\Leb^2(0,1;H_\Omega)$
  and a real $U\in\C([0,1];V_\Omega)\cap\C^1([0,1];V_\Omega')\cap\W^{1,1}(0,1;H_\Omega)$, in
  particular in $\Leb^2(0,1;V_\Omega)\cap\Sob^1(0,1;V_\Omega')$, with $U(0)=0$, which is the unique
  Lions variational solution of $U'(t)-\Div(\frB(t)\nabla U(t))=F(t)$ in $V_\Omega'$ for every
  $t\in[0,1]$, but $U'\notin\Leb^2(0,1;H_\Omega)$ and
  $-\Div(\frB\nabla U)\notin\Leb^2(0,1;H_\Omega)$.
\end{corollary}

\textbf{Lean form.}  The statement is that of
\texttt{LionsC12Challenge.lions\_counterexample\_domain} in \texttt{Challenge.lean}, proved in
\texttt{Solution.lean} by \texttt{LionsC12.lions\_counterexample\_domain}.  ``Bounded domain'' is a bounded non-empty open set; neither boundedness nor
connectedness is used in the proof (\texttt{lions\_counterexample\_open}).  The matrix norm in $\Leb^\infty(\Omega;\R^{d\times d})$ is the operator norm of $\R^d$
(every norm on $\R^{d\times d}$ gives the same space; for $\frb\Id_d-\frb'\Id_d$ the operator norm
is $\abs{\frb-\frb'}$).  The clauses are read as in Theorem~\ref{thm:main}.

\begin{proof}
  \leanok
  \uses{prop:line, prop:tensor, lem:domain-clauses, prop:holder-sharp,
    lem:solution-regularity, lem:lions-uniqueness-domain}
  Let $W$, $G$ be the solution and forcing on $\R^d$ of Proposition~\ref{prop:line} ($d=1$) or
  Proposition~\ref{prop:tensor} ($d\geq2$), supported in a fixed compact set $K$ in space.
  Choose $x_0$ and $0<r\leq1$ with $x_0+rK\Subset\Omega$, and put $s=r^{-2}t$,
  $z=r^{-1}(x-x_0)$, $\frb(t,x)=\fra(s,z_1)$, $U(t,x)=W(s,z)$, $F(t,x)=r^{-2}G(s,z)$.  The
  parabolic change of variables gives the equation; the support in $x_0+rK$ gives the
  regularity in $V_\Omega$, $H_\Omega$ and $V_\Omega'$; the time rescaling gives
  $\norm{\frB-\Id_d}_{\C^{0,1/2}}\leq Cr^{-1}\eps$, so $\eps$ is chosen small in terms of $r$ and
  $\delta$; the failure of $\alpha$-H\"older continuity and
  $\int_0^1\norm{U'}_{H_\Omega}^2=r^{d-2}\int_0^1\norm{\partial_sW}^2=\infty$ follow from
  Proposition~\ref{prop:holder-sharp} and Lemma~\ref{lem:solution-regularity}.  Uniqueness is
  Lemma~\ref{lem:lions-uniqueness-domain}.
\end{proof}


\chapter{Spatial oscillatory modes}
The construction uses the profile $\phi(x)=\sin^2x$ and, for integers $K>2$, the modes
\[
  p_K(x)=\sin(Kx)\sin(2x)=\sin(Kx)\,\partial_x\phi(x),\qquad c_K(x)=\cos(Kx)\cos(2x).
\]
$\phi$, $p_K$ and their first derivatives vanish at $0$ and $\pi$
(\texttt{eq:endpoint-conditions} of the paper).

\begin{definition}[The profile and the modes]
  \label{def:modes}
  \lean{LionsC12.φ, LionsC12.φ', LionsC12.p, LionsC12.p', LionsC12.p'', LionsC12.c,
    LionsC12.pV, LionsC12.φV}
  \leanok
  \uses{def:spaces}
  The real functions $\phi$, $p_K$, $c_K$ and their derivatives, and $\phi$, $p_K$ as
  elements of $V$ (they are $\C^1$ on $\R$ and vanish at $0$ and $\pi$; the zero extension
  outside $[0,\pi]$ is $\C^1$, and a dilation towards the midpoint approximates it by
  functions in $\C_c^1(0,\pi)$).
\end{definition}

\begin{lemma}[Oscillatory-mode identities]
  \label{lem:oscillatory-modes}
  \lean{LionsC12.lemma_2_1, LionsC12.normSq_p, LionsC12.normSq_c, LionsC12.normSq_p',
    LionsC12.norm_p_Hm1_le, LionsC12.isRealisation_pV, LionsC12.inner_p_p,
    LionsC12.inner_p_c', LionsC12.inner_c_p, LionsC12.inner_c_c}
  \leanok
  \uses{def:modes}
  There is an absolute constant $C>0$ such that, for every integer $K>2$,
  \[
    \norm{p_K}_H^2=\norm{c_K}_H^2=\frac\pi4,\qquad
    \norm{\partial_xp_K}_H^2=\frac\pi4(K^2+4),\qquad \norm{p_K}_{V'}\leq CK^{-1}.
  \]
  Moreover $A_{0,H}p_K=(K^2+4)p_K-4Kc_K$.  If $L\geq16K$, then
  $\mathrm{span}\{p_K,c_K\}$ and $\mathrm{span}\{p_L,c_L\}$ are orthogonal in $H$.
\end{lemma}

\textbf{Lean form.}  The constant is $C=6$; the norm of $V'$ is the operator norm of the
anti-dual of $V$ with the graph norm.  $A_{0,H}p_K=g$ is the statement that $p_K\in V$ and
$\langle\partial_xp_K,\partial_xw\rangle_H=\langle g,w\rangle_H$ for every $w\in V$
(\texttt{IsRealisation one 0}); this is the realisation in $H$ of the form operator.

\begin{proof}
  \leanok
  Product-to-sum formulas reduce the norms and the orthogonality to the cosine integrals on
  $(0,\pi)$.  For the dual norm, $P_K=\frac12\bigl(\frac{\sin((K-2)x)}{K-2}-\frac{\sin((K+2)x)}{K+2}\bigr)$
  satisfies $P_K'=p_K$ and $P_K(0)=P_K(\pi)=0$, so
  $\abs{\langle p_K,w\rangle}=\abs{\langle P_K,\partial_xw\rangle}\leq\norm{P_K}_H\norm{w}_V$ for
  real $\C^1$ test functions, and the complex test functions and the closure are handled by real
  and imaginary parts and continuity.  The formula for $A_{0,H}p_K$ is the differentiation of
  $\sin(Kx)\sin(2x)$ and integration by parts in $V$.
\end{proof}

\begin{lemma}[The coefficient acting on the profile]
  \label{lem:coefficient-profile}
  \lean{LionsC12.hasDerivAt_coeff_profile}
  \leanok
  \uses{def:modes}
  $L_{K^{-1}\cos(Kx)}\phi=p_K-2K^{-1}c_K$
  (\texttt{eq:coefficient-profile-identity} of the paper).
\end{lemma}

\begin{proof}
  \leanok
  Differentiate $K^{-1}\cos(Kx)\,\partial_x\phi(x)$.
\end{proof}


\chapter{Time blocks, frequencies, and the coefficient}
For $j\geq1$ let $\tau_j=\frac5{16}-\frac1{16j}$, $\ell_j=\tau_{j+1}-\tau_j=\frac1{16j(j+1)}$,
$M_j=j$, $I_j=(\tau_j,\tau_{j+1})$, $J_j=[\tau_j+\ell_j/3,\tau_j+2\ell_j/3]$, and for
$1\leq m\leq M_j$ let $K_{j,m}=16^{m+1}j(j+1)$ and $\theta_{j,m}(t)=K_{j,m}^2(t-\tau_j)$.  With a
cutoff $\eta\in\C_c^\infty(0,1)$, $0\leq\eta\leq1$, $\supp\eta\subseteq[1/10,9/10]$,
$\eta=1$ on $[1/3,2/3]$, and $\eta_j(t)=\eta((t-\tau_j)/\ell_j)$,
\[
  \beta(t,x)=\sum_{j\geq1}\sum_{m=1}^{M_j}\eta_j(t)K_{j,m}^{-1}\sin\theta_{j,m}(t)\cos(K_{j,m}x),
  \qquad \fra(t,x)=1+\eps\beta(t,x).
\]

\begin{definition}[Time blocks, frequencies, cutoffs and the coefficient]
  \label{def:coefficient}
  \lean{LionsC12.τ, LionsC12.ℓ, LionsC12.M, LionsC12.K, LionsC12.θ, LionsC12.I, LionsC12.J,
    LionsC12.S, LionsC12.Cutoffs, LionsC12.Cutoffs.nonempty, LionsC12.ηj, LionsC12.βmode,
    LionsC12.βblock, LionsC12.β, LionsC12.βx, LionsC12.a}
  \leanok
  The objects above.  At a fixed time at most one block is active, so the sum is pointwise
  finite.
\end{definition}

\textbf{Lean form.}  The two cutoffs $\eta$ and $\rho$ of the paper are recorded, with the
properties used in the proofs, as the structure \texttt{Cutoffs}; its existence
(\texttt{Cutoffs.nonempty}) is proved from Mathlib's smooth bump functions, and the theorems
are proved for every inhabitant.  The infinite sums over blocks are \texttt{finsum}s.

\begin{lemma}[Arithmetic of the blocks]
  \label{lem:reciprocal-sums}
  \lean{LionsC12.S_le, LionsC12.sum_K_inv_sq_le, LionsC12.M_mul_ℓ, LionsC12.M_mul_S_le,
    LionsC12.K_succ, LionsC12.one_le_K_sq_mul_ℓ, LionsC12.I_disjoint}
  \leanok
  \uses{def:coefficient}
  $S_j=\sum_mK_{j,m}^{-1}\leq\ell_j/15$, $\sum_mK_{j,m}^{-2}\leq\ell_j^2/255$,
  $M_j\ell_j=\frac1{16(j+1)}$ and $M_jS_j\leq\frac1{240(j+1)}$
  (\texttt{eq:reciprocal-sums}); consecutive frequencies differ by the factor $16$; the
  blocks are disjoint.
\end{lemma}

\begin{proof}
  \leanok
  Geometric sums.
\end{proof}

The paper's Remark~3.1 (\emph{Choice of the frequencies}) observes that the factor $16$ is used
only for summability and separation, and that a dyadic family would do after changing
constants.  It makes no claim used later and has no Lean counterpart; the formalisation uses the
factor $16$ of the paper.

\begin{proposition}[Ellipticity and $\frac12$-H\"older time regularity]
  \label{prop:coefficient}
  \lean{LionsC12.abs_β_le, LionsC12.abs_a_sub_one_le, LionsC12.abs_β_sub_le,
    LionsC12.abs_a_sub_le, LionsC12.eLpNorm_a_sub_one_le, LionsC12.eLpNorm_a_sub_le,
    LionsC12.isBoundedCoeff_a, LionsC12.a_pos}
  \leanok
  \uses{def:coefficient, lem:reciprocal-sums}
  $\fra$ is bounded and real, $\norm{\fra-1}_{\Leb^\infty(\R^2)}\leq C\eps$, and
  $\norm{\fra(t,\cdot)-\fra(s,\cdot)}_{\Leb^\infty(\R)}\leq C\eps\abs{t-s}^{1/2}$ for $s,t\in\R$,
  with $C$ depending only on $\eta$.  For small $\eps$, $\fra$ is uniformly elliptic.
\end{proposition}

\textbf{Lean form.}  The bounds are pointwise for every $(t,x)$, with the explicit constants
$\abs{\fra-1}\leq\eps/480$ and $\abs{\fra(t,x)-\fra(s,x)}\leq2\eps C_\beta\abs{t-s}^{1/2}$, where
$C_\beta$ depends only on the cutoff; the $\Leb^\infty(0,\pi)$ forms
(\texttt{eLpNorm\_a\_sub\_one\_le}, \texttt{eLpNorm\_a\_sub\_le}) follow.  Uniform ellipticity
holds for every $0\leq\eps\leq1$.

\begin{proof}
  \leanok
  At each time one block contributes, so $\abs{\beta}\leq\sup_jS_j\leq\ell_1/15$.  For the
  H\"older bound, compare $\abs{t-s}$ with $K_{j,m}^{-2}$: for the modes with
  $K_{j,m}^2\abs{t-s}\leq1$ use the derivative bound, for the others the uniform bound, and sum
  the two geometric series; times in different blocks are reduced to the block endpoints.
\end{proof}

\begin{proposition}[Sharpness within the H\"older scale]
  \label{prop:holder-sharp}
  \lean{LionsC12.not_holder, LionsC12.not_holder_Linfty}
  \leanok
  \uses{def:coefficient}
  For every $\frac12<\alpha<1$, $\fra\notin\C^{0,\alpha}([0,1];\Leb^\infty(0,\pi))$.
\end{proposition}

\textbf{Lean form.}  \texttt{not\_holder} shows that the pointwise quotient
$\abs{\fra(t,0)-\fra(s,0)}/\abs{t-s}^\alpha$ is unbounded on $[0,1]$, at the point $x=0$, with
$t,s$ in the middle third of a block; \texttt{not\_holder\_Linfty} deduces the statement for the
essential supremum on $(0,\pi)$, because $\fra(t,\cdot)$ is continuous and the essential supremum
of a continuous function on $(0,\pi)$ dominates its value at $0$.  The paper argues with the
Fourier coefficient $\Lambda_K$; the point evaluation at $x=0$ is a different but equivalent route,
since all modes of a block are evaluated at $\cos(K\cdot0)=1$.

\begin{proof}
  \leanok
  Inside $J_j$ only the first mode of block $j$ moves at the scale $h=\pi/(2K_{j,1}^2)$; the
  increment over $h$ is at least $c\eps K_{j,1}^{-1}=c'\eps h^{1/2}$, while $h^\alpha\ll h^{1/2}$.
\end{proof}

\begin{proposition}[Common operator domain without uniform graph control]
  \label{prop:frozen-domains}
  \lean{LionsC12.IsRealisation, LionsC12.domain_eq, LionsC12.MemH2, LionsC12.domain_one_eq_H2,
    LionsC12.domain_eq_H2, LionsC12.exists_bound_on_D₀, LionsC12.not_uniform_bound_on_D₀}
  \leanok
  \uses{def:coefficient, def:lions-solution, lem:coefficient-profile, prop:sobolev}
  Let $D_0=D(A_{0,H})=\Sob^2(0,\pi)\cap\Sob^1_0(0,\pi)$ with the graph norm of $A_{0,H}$.  Then
  $D(A_H(t))=D_0$ for $t\in[0,1]$, but $\sup_{t\in[0,1]}\norm{A_H(t)}_{\cL(D_0,H)}=\infty$.
\end{proposition}

\textbf{Lean form.}  $D(A_H(t))$ is the set of $v\in V$ for which some $h\in H$ satisfies
$\fra(t)[v,w]=\langle h,w\rangle_H$ for all $w\in V$, and $D_0=D(A_{0,H})$ is the same set for
$\fra=1$ (\texttt{IsRealisation}), and \texttt{domain\_eq} proves $D(A_H(t))=D_0$.
\texttt{domain\_one\_eq\_H2} identifies $D_0$ with $\Sob^2\cap\Sob^1_0$, the definition in the
paper, where $v\in\Sob^2$ means that $v$ and its weak derivative have weak derivatives in $\Leb^2$
(\texttt{MemH2}, Proposition~\ref{prop:sobolev}), and \texttt{domain\_eq\_H2} states
$D(A_H(t))=\Sob^2\cap\Sob^1_0$ for every $t$.
$\norm{A_H(t)}_{\cL(D_0,H)}<\infty$ is \texttt{exists\_bound\_on\_D₀}:
$\norm{A_H(t)v}_H\leq C_t(\norm{v}_V+\norm{A_{0,H}v}_H)$, and \texttt{not\_uniform\_bound\_on\_D₀}
says that no $C$ works for all $t\in[0,1]$.  The graph norm is taken with $\norm{v}_V$ in place of
$\norm{v}_H$, an equivalent norm on $D_0$.

\begin{proof}
  \leanok
  For fixed $t$, $\fra_t,\fra_t^{-1}\in\W^{1,\infty}$, and $A_H(t)v=-\fra_t\partial_x^2v-(\partial_x\fra_t)\partial_xv$,
  $A_{0,H}v=\fra_t^{-1}(A_H(t)v+(\partial_x\fra_t)\partial_xv)$ in the weak sense, tested against
  $\fra_t^{\pm1}w$.  For the blow-up, $A_H(t)\phi$ contains
  $\eps\sum_m\eta_j\sin\theta_{j,m}\,p_{K_{j,m}}$ by Lemma~\ref{lem:coefficient-profile}, whose
  norm in $H$ grows like $\eps\sqrt{M_j}$ on $J_j$ by orthogonality.
\end{proof}


\chapter{The variational solution and its time derivative}
With a second cutoff $\rho\in\C_c^\infty(0,1)$, $\rho=1$ on $[1/4,1/2]$, and
\[
  y_{j,m}(t)=-\frac{\eps}{2K_{j,m}^2}\eta_j(t)\bigl(\sin\theta_{j,m}(t)-\cos\theta_{j,m}(t)\bigr),
  \qquad v=\sum_{j,m}y_{j,m}\,p_{K_{j,m}},\qquad u=\rho\phi+v,
\]
the functions $y_{j,m}$ solve $\partial_ty_{j,m}+K_{j,m}^2y_{j,m}=-\eps\eta_j\sin\theta_{j,m}+r_{j,m}$
with $\abs{r_{j,m}}\leq C\eps K_{j,m}^{-2}\ell_j^{-1}$ (\texttt{eq:y-ode}, \texttt{eq:r-def}).

\begin{definition}[The candidate solution]
  \label{def:solution}
  \lean{LionsC12.y, LionsC12.y', LionsC12.r, LionsC12.vblock, LionsC12.v, LionsC12.u,
    LionsC12.ut, LionsC12.ux, LionsC12.uxx, LionsC12.uV, LionsC12.uV', LionsC12.utH}
  \leanok
  \uses{def:coefficient, def:modes}
  The real functions $y_{j,m}$, $v$, $u$ with their classical derivatives in $t$ and $x$; $u(t)$ as
  an element of $V$ (\texttt{uV}), the derivative of $t\mapsto u(t)\in V'$ (\texttt{uV'}) and
  the time derivative $\partial_tu(t)\in H$ (\texttt{utH}).
\end{definition}

\begin{lemma}[Regularity and failure of $\Leb^2$-integrability of the time derivative]
  \label{lem:solution-regularity}
  \lean{LionsC12.continuous_uV, LionsC12.hasDerivAt_ofH10_uV, LionsC12.continuous_uV',
    LionsC12.integrableOn_utH, LionsC12.hasWeakDeriv_val_uV, LionsC12.uV_zero,
    LionsC12.not_integrableOn_utNormSq, LionsC12.not_memLp_utH,
    LionsC12.not_exists_L2_rep_uV'}
  \leanok
  \uses{def:solution, lem:oscillatory-modes, lem:reciprocal-sums}
  $u\in\C([0,1];V)\cap\C^1([0,1];V')\cap\W^{1,1}(0,1;H)$, $u(0)=0$, and
  $\partial_tu\notin\Leb^2(0,1;H)$.
\end{lemma}

\textbf{Lean form.}  $u$ is continuous into $V$ on all of $\R$, and $t\mapsto u(t)\in V'$ has the
derivative \texttt{uV'} at every $t\in\R$, continuous in $t$.  $\W^{1,1}(0,1;H)$ is
Bochner integrability of $t\mapsto u(t)\in H$ and of $\partial_tu$, together with
$\partial_tu$ being the $H$-valued distributional derivative.  ``$\partial_tu\notin\Leb^2$'' is
\texttt{not\_exists\_L2\_rep\_uV'}: no $g\in\Leb^2(0,1;H)$ represents $u'$ almost everywhere.

\begin{proof}
  \leanok
  On each block the sum is finite and smooth, $\abs{y_{j,m}}\leq C\eps K_{j,m}^{-2}$ and
  $\abs{\partial_ty_{j,m}}\leq C\eps$, so by Lemma~\ref{lem:oscillatory-modes}
  $\norm{v_j}_{\C(\R;V)}+\norm{\partial_tv_j}_{\C(\R;V')}\leq C\eps\sum_mK_{j,m}^{-1}\leq C\eps\ell_j$;
  the series converge uniformly in $V$ and $V'$, which gives $u\in\C([0,1];V)\cap\C^1([0,1];V')$.
  In $H$, $\norm{\partial_tv_j(t)}_H^2\leq C\eps^2M_j$ by orthogonality, so
  $\norm{\partial_tv_j}_{\Leb^1(I_j;H)}\leq C\eps\ell_j\sqrt{M_j}$, which is summable.  On $J_j$,
  $\partial_tu=-\frac\eps2\sum_m(\cos\theta_{j,m}+\sin\theta_{j,m})p_{K_{j,m}}$ and
  $\int_{J_j}(\cos\theta_{j,m}+\sin\theta_{j,m})^2\geq\ell_j/6$, so
  $\int_0^1\norm{\partial_tu}_H^2\geq\frac{\pi\eps^2}{96}\sum_jM_j\ell_j=\infty$.
\end{proof}


\chapter{Cancellation and the forcing}
The forcing is defined by $f=\partial_tu+\cA(\cdot)u$.  The point of the construction is the exact
cancellation of each mode against the coefficient acting on the profile,
\[
  (\partial_t+A_{0,H})(y_{j,m}p_K)+\eps L_{\eta_jK^{-1}\sin\theta_{j,m}\cos(Kx)}(\rho\phi)=E_{j,m},
\]
with a remainder $E_{j,m}$ that is small in $H$ (\texttt{eq:exact-mode-cancellation}).

\begin{lemma}[Exact cancellation and the equation pointwise]
  \label{lem:cancellation}
  \lean{LionsC12.mode_cancellation, LionsC12.pde, LionsC12.f, LionsC12.f0, LionsC12.E,
    LionsC12.R, LionsC12.Lβv, LionsC12.Au}
  \leanok
  \uses{def:solution, lem:coefficient-profile}
  The identity above holds for every mode, and summing gives
  $f=f_0+R+\eps L_{\beta}v$ (\texttt{eq:f-decomposition}); pointwise,
  $\partial_tu(t,x)-\partial_x(\fra\partial_xu)(t,x)=f(t,x)$ for every $(t,x)$.
\end{lemma}

\begin{proof}
  \leanok
  Lemma~\ref{lem:oscillatory-modes}, Lemma~\ref{lem:coefficient-profile} and the ODE for $y_{j,m}$.
\end{proof}

\begin{lemma}[$u(t)$ lies in the domain of $A_H(t)$]
  \label{lem:domain-membership}
  \lean{LionsC12.AuH, LionsC12.form_uV, LionsC12.integrableOn_AuH, LionsC12.not_memLp_AuH,
    LionsC12.AH_u_L1_not_L2, LionsC12.not_integrableOn_normSq_Au}
  \leanok
  \uses{def:solution, def:lions-solution, prop:forcing, lem:solution-regularity}
  For every $t$, $\fra(t)[u(t),w]=\langle-\partial_x(\fra\partial_xu)(t),w\rangle_H$ for all $w\in V$, so
  $u(t)\in D(A_H(t))$; and $A_H(\cdot)u(\cdot)\in\Leb^1(0,1;H)\setminus\Leb^2(0,1;H)$.
\end{lemma}

\begin{proof}
  \leanok
  Integration by parts in $V$ against the $\C^1$ flux $\fra\partial_xu(t)$.  Then
  $A_Hu=f-\partial_tu$ with $f\in\C([0,1];H)$ (Proposition~\ref{prop:forcing}) and
  $\partial_tu\in\Leb^1\setminus\Leb^2$ (Lemma~\ref{lem:solution-regularity}).
\end{proof}

\begin{proposition}[Continuous $H$-valued forcing]
  \label{prop:forcing}
  \lean{LionsC12.fH, LionsC12.continuous_fH, LionsC12.im_fH, LionsC12.utH_add_AuH,
    LionsC12.equation_uV, LionsC12.isLionsSolution_uV, LionsC12.eq_uV_of_isLionsSolution}
  \leanok
  \uses{lem:cancellation, lem:solution-regularity, prop:coefficient, lem:lions-uniqueness}
  $f\in\C([0,1];H)$, and $\partial_tu(t)+\cA(t)u(t)=f(t)$ in $V'$ for every $t\in[0,1]$.
\end{proposition}

\textbf{Lean form.}  $f$ is jointly continuous on $\R^2$ and $\C([0,1];H)$ is deduced from it;
the equation in $V'$ holds for every $t\in\R$.  In addition, \texttt{isLionsSolution\_uV} places
$u$ in the Lions class and \texttt{eq\_uV\_of\_isLionsSolution} records uniqueness.

\begin{proof}
  \leanok
  $\norm{R_j(t)}_H\leq C\eps\ell_j$ and the supports are disjoint, so $R$ is continuous into $H$.
  On block $j$, $\norm{\beta}_\infty\leq S_j$, $\norm{\partial_x\beta}_\infty\leq M_j$,
  $\norm{\partial_xv}_H\leq C\eps S_j$ and $\norm{\partial_x^2v}_H\leq C\eps M_j$, so
  $\norm{\eps L_\beta v}_H\leq C\eps^2M_jS_j\leq C\eps^2/(j+1)$, and the block terms, continuous
  with disjoint supports, sum to an element of $\C([0,1];H)$.  The equation in $V'$ is
  Lemma~\ref{lem:cancellation} tested against $w\in V$, by Lemma~\ref{lem:domain-membership}.
\end{proof}

\begin{remark}[Growth of the maximal-regularity constant]
  \label{rem:growth}
  \lean{LionsC12.aN, LionsC12.uN, LionsC12.uNV, LionsC12.fNH, LionsC12.IsMaxRegConstant,
    LionsC12.maxRegConstant_growth, LionsC12.maximal_regularity_constant_growth,
    LionsC12.integral_normSq_fN_le, LionsC12.truncation_maximal_regularity,
    LionsC12.truncation_maxReg_growth, LionsC12.truncation_optimalMaxRegConstant,
    LionsC12.isMaxRegConstant_of_C1}
  \leanok
  \uses{prop:coefficient, lem:solution-regularity, prop:forcing, prop:frozen-domains,
    thm:lions-garding, lem:lions-uniqueness}
  For $N\in\N$ let $\fra_N=1+\eps\sum_{j\leq N}\beta_j$, $u_N=\rho\phi+\sum_{j\leq N}v_j$ and
  $f_N=\partial_tu_N+\cA_N(t)u_N$.  The problem associated with $\cA_N$ has maximal
  $\Leb^2$-regularity.  Moreover
  $\sup_N\norm{\fra_N-1}_{\C^{0,1/2}([0,1];\Leb^\infty)}\leq C\eps$,
  $\sup_N\norm{f_N}_{\Leb^2(0,1;H)}<\infty$, and
  $\norm{\partial_tu_N}^2_{\Leb^2(0,1;H)}\geq c\eps^2\log(N+1)$; so every maximal-regularity
  constant $C_N$ of $\fra_N$ satisfies $C_N\geq c\eps\sqrt{\log(N+1)}$.
\end{remark}

\textbf{Lean form.}  The explicit constants are $c=\pi/1536$ and
$\norm{\partial_tu_N}^2_{\Leb^2(0,1;H)}\geq\frac{\pi\eps^2}{1536}(\log(N+2)-1)$.  A
maximal-regularity constant of $\fra_N$ (\texttt{IsMaxRegConstant}) is a $C$ with
$\norm{w'}_{\Leb^2(0,1;H)}+\norm{\cA_N w}_{\Leb^2(0,1;H)}\leq C\norm{g}_{\Leb^2(0,1;H)}$ for every
Lions solution $w$ of $w'+\cA_Nw=g$, $w(0)=0$, whose derivative and elliptic part lie in
$\Leb^2(0,1;H)$.  Maximal regularity of $\fra_N$ (\texttt{truncation\_maximal\_regularity}) is:
for every $g\in\Leb^2(0,1;H)$ there is a Lions solution with derivative and elliptic part in
$\Leb^2(0,1;H)$ and $\norm{w'}+\norm{\cA_Nw}\leq C\norm{g}$, and \texttt{IsMaxRegConstant} holds
for that finite $C$; \texttt{truncation\_maxReg\_growth} combines it with the lower bound, and the
best constant is attained (\texttt{truncation\_optimalMaxRegConstant}).  The paper argues from
smoothness of the finite sums; Lean uses only that $\fra_N(t,\cdot)$ is $\C^1$ with
$\abs{\partial_x\fra_N}\leq\eps\sum_{j\leq N}j$ uniformly in $t$, and applies Lions' theorem in the
triple $D_0\subset V\subset D_0'$ with pivot $V$, $D_0=\Sob^2\cap\Sob^1_0$
(Theorem~\ref{thm:lions-garding}): there $\re\langle\cA_N(t)u,(I+A_{0,H})u\rangle_H$ satisfies a
Gårding inequality in the graph norm of $A_{0,H}$, the solution has its time derivative in
$\Leb^2(0,1;H)$ because $I+A_{0,H}:D_0\to H$ is onto, and uniqueness of the Lions solution transfers
the regularity to every solution.  No time regularity of $\fra_N$ is used; for the full
coefficient the argument fails because $\sup_t\norm{\partial_x\fra(t)}_\infty=\infty$.  The constant
is explicit but far from optimal.

\begin{proof}
  \leanok
  Maximal regularity as described above.  For the lower bound, the bounds of
  Proposition~\ref{prop:coefficient} and Proposition~\ref{prop:forcing} are uniform in $N$, and
  the lower bound of Lemma~\ref{lem:solution-regularity} on the first $N$ blocks gives
  $\sum_{j\leq N}\ell_jM_j\geq c\log(N+1)$.
\end{proof}


\chapter{Full-space and bounded-domain consequences}
The interval counterexample is extended by zero to the line, tensorised with a bump in the
remaining variables, and rescaled into a ball inside an arbitrary open set.  The key ingredient is
that $\phi$, $p_K$ and their first derivatives vanish at both endpoints, so that the zero extension
of the flux $\fra\partial_xu$ has no Dirac masses at $0$ and $\pi$.  For a function $z$ on $(0,\pi)$,
$E_0z$ is its zero extension to $\R$.

\begin{lemma}[Zero extension of a vanishing flux]
  \label{lem:zero-extension-flux}
  \lean{LionsC12.zero_extension_flux, LionsC12.integral_deriv_mul_primitive}
  \leanok
  Let $q\in\Leb^1(0,1;\Sob^1(0,\pi))$ with $q(t,0)=q(t,\pi)=0$ for almost every $t\in(0,1)$.  Then
  $\partial_x(E_0q)=E_0(\partial_xq)$ in $\mathcal D'((0,1)\times\R)$.
\end{lemma}

\textbf{Lean form.}  $q\in\Leb^1(0,1;\Sob^1(0,\pi))$ is read as: $q$ and $\partial_xq$ are integrable
on $(0,1)\times(0,\pi)$, and for almost every $t$, $q(t,\cdot)$ is the primitive
$x\mapsto\int_0^x\partial_xq(t,y)\dd y$ on $[0,\pi]$ (the absolutely continuous representative with
$q(t,0)=0$) with $q(t,\pi)=0$; this is $\Leb^1(0,1;\W^{1,1}(0,\pi))$, which contains
$\Leb^1(0,1;\Sob^1(0,\pi))$.  The identity is
$\int_0^1\!\int_\R E_0q\,\partial_x\psi=-\int_0^1\!\int_\R E_0(\partial_xq)\,\psi$ for every $\psi$ that is
$\C^1$ in $x$, which contains the test functions of $\mathcal D'((0,1)\times\R)$.

\begin{proof}
  \leanok
  Fubini on the triangle $\{0<y<x<\pi\}$: $\int_0^\pi\psi'(x)\int_0^xg=\psi(\pi)\int_0^\pi g-\int_0^\pi\psi g$,
  and the boundary term vanishes because $q(t,\pi)=0$.
\end{proof}

\begin{lemma}[The one-dimensional weak form on the line]
  \label{lem:weak-line}
  \lean{LionsC12.weak1D, LionsC12.uZ, LionsC12.fZ}
  \leanok
  \uses{lem:zero-extension-flux, lem:cancellation}
  For every $t$, the zero extensions $\widetilde u=E_0u$ and $\widetilde f=E_0f$ satisfy
  $\int_\R\fra\,\partial_x\widetilde u\,\overline{\psi'}=\int_\R(\widetilde f-\partial_t\widetilde u)\overline\psi$ for every
  $\psi\in\C^1(\R)$ with no boundary condition, also after the parabolic rescaling
  $(t,x)\mapsto(t/r^2,(x-x_0)/r)$.
\end{lemma}

\begin{proof}
  \leanok
  $\widetilde u(t,\cdot)$ is $\C^1$ because $u$ and $\partial_xu$ vanish at $0$ and $\pi$, and the flux
  $\fra\partial_x\widetilde u$ is $\C^1$ on $[0,\pi]$ and vanishes at the endpoints; integrate by parts
  on $[0,\pi]$ as in Lemma~\ref{lem:zero-extension-flux}.
\end{proof}

\begin{lemma}[Tensor products]
  \label{lem:tensor}
  \lean{LionsC12.tens, LionsC12.lapQ, LionsC12.tensorΩ_weak}
  \leanok
  \uses{lem:weak-line}
  For bounded measurable $P,P',A,b$ on $\R$ with the one-dimensional weak identity
  $-\partial_1(b\partial_1P)=A$ and $Q\in\C_c^\infty(\R^{d-1})$,
  $-\Div\bigl(b(x_1)\nabla(P\otimes Q)\bigr)=A\otimes Q-bP\otimes\Delta'Q$ weakly on $\R^d$.
\end{lemma}

\begin{proof}
  \leanok
  Fubini, the identity in $x_1$, and integration by parts on lines in the $x'$ variables.
\end{proof}

\begin{definition}[The construction on $\R^d$ and its rescaling]
  \label{def:domain-construction}
  \lean{LionsC12.bumpE, LionsC12.domB, LionsC12.domW, LionsC12.domWt, LionsC12.domGrad,
    LionsC12.domF, LionsC12.DomParams, LionsC12.dU, LionsC12.dG, LionsC12.dF, LionsC12.dGv,
    LionsC12.dU'}
  \leanok
  \uses{def:spaces-domain, def:solution}
  For $d=n+1$, a bump $\chi$ on $\R^n$, a centre $p$ and a radius $0<r\leq1$:
  $\frb(t,x)=\fra(t/r^2,(x_1-p_1)/r)$, $U(t,x)=\widetilde u(t/r^2,(x_1-p_1)/r)\,\chi((x'-p')/r)$, and
  $F=r^{-2}\bigl(\widetilde f\otimes\chi-\fra\widetilde u\otimes\Delta'\chi\bigr)$ in the rescaled variables, as
  functions on $\R\times\R^d$ and as elements of $V_\Omega$, $H_\Omega$, $V_\Omega'$.
\end{definition}

\begin{proposition}[Counterexample on the real line]
  \label{prop:line}
  \lean{LionsC12.lions_counterexample_fullspace}
  \leanok
  \uses{def:domain-construction, lem:weak-line, lem:solution-regularity, prop:forcing,
    lem:lions-uniqueness-domain}
  $\widetilde u=E_0u$ is the Lions variational solution on $\R$ with forcing $\widetilde f=E_0f$; $\fra$ is
  bounded, uniformly elliptic and in $\C^{0,1/2}([0,1];\Leb^\infty(\R))$, but
  $\partial_t\widetilde u\notin\Leb^2(0,1;\Leb^2(\R))$.
\end{proposition}

\begin{proposition}[Counterexamples in full space]
  \label{prop:tensor}
  \lean{LionsC12.lions_counterexample_fullspace, LionsC12.DomainClauses}
  \leanok
  \uses{def:domain-construction, lem:tensor, lem:solution-regularity, prop:forcing,
    lem:lions-uniqueness-domain}
  Let $d\geq2$.  $\frB_d(t,x)=\fra(t,x_1)\Id_d\in\C^{0,1/2}([0,1];\Leb^\infty(\R^d;\R^{d\times d}))$ is
  real symmetric, bounded and uniformly elliptic, and there are a real
  $F\in\Leb^2(0,1;\Leb^2(\R^d))$ and a unique Lions solution $U$ with $U(0)=0$ of
  $\partial_tU-\Div(\frB_d\nabla U)=F$, with $\partial_tU\notin\Leb^2(0,1;\Leb^2(\R^d))$.
\end{proposition}

\textbf{Lean form.}  One theorem, \texttt{lions\_counterexample\_fullspace}, proves both
propositions for every $d=n+1\geq1$ on $\Omega=\R^d$ without rescaling: the coefficient is
$\frb(t,x)=\fra(t,x_1)$, $U(t)=\widetilde u(t)\otimes\chi$ and
$F(t)=\widetilde f(t)\otimes\chi-\fra\widetilde u(t)\otimes\Delta'\chi$, and they satisfy every clause of
Corollary~\ref{cor:bounded-domains} on $\R^d$ (\texttt{DomainClauses}), which is more than the
propositions assert.  The case $d=1$ is $n=0$: $\R^0$ is a point, $\chi\equiv1$ and $\Delta'\chi=0$,
so $U=\widetilde u$ and $F=\widetilde f$; the line is \texttt{EuclideanSpace ℝ (Fin 1)} rather than
$\R$.  The bump is a smooth bump of Mathlib supported in the closed unit ball.

\begin{proof}
  \leanok
  The spatial weak formulation at fixed $t$ follows from Lemmas~\ref{lem:weak-line}
  and~\ref{lem:tensor}, for test functions first and for all of $V_\Omega$ by density
  (\texttt{swf}).  The regularity of $U$ and $F$ follows from that of $\widetilde u$ and
  $\widetilde f$, because every function involved is supported in a fixed compact set.  The negative
  clause is $\int\norm{\partial_tU}^2=\norm{\chi}^2_{\Leb^2}\int\norm{\partial_t\widetilde u}^2=\infty$.
\end{proof}

\textbf{Lean form of the proof of Corollary~\ref{cor:bounded-domains}.}  The statement is proved for
every non-empty open $\Omega\subseteq\R^d$ (\texttt{lions\_counterexample\_open}); boundedness and
connectedness are not used.  The point $x_0=p$ is the centre of a ball $B(p,\rho)\subseteq\Omega$,
$r=\min\{1,\rho/(2(\pi+1))\}$, and the compact support set is $\overline B(p,r(\pi+1))$ in place of
$x_0+rK$.  $\eps=\min\{1,\delta r/(2(1+2C_\beta))\}$, the H\"older constant is $B=2\eps C_\beta/r$ and
$A=\eps/480$.  The failure of $\alpha$-H\"older continuity evaluates $\frb$ at $p$, where the
rescaled first coordinate is $0$ (as in Proposition~\ref{prop:holder-sharp}).  For
$U\in\C^1([0,1];V_\Omega')$ the paper uses $W\in\C(\R;\Sob^1)$, $\partial_sW\in\C(\R;\Sob^{-1})$ and
bounded dilations; Lean instead defines $U'(t)=F(t)-\cA(t)U(t)$, which is continuous into
$V_\Omega'$ because $t\mapsto\frb\nabla U(t)$ is continuous in $\Leb^2(\Omega;\CC^d)$, identifies it
with $\partial_tU(t)\in H_\Omega$ by the spatial weak formulation (\texttt{dU'\_eq}), and obtains the
derivative at the accumulation time $5r^2/16$ from the fundamental theorem of calculus
(\texttt{hasDerivAt\_ofH10\_dU}).

\begin{lemma}[The clauses of Corollary~1.2 for the construction]
  \label{lem:domain-clauses}
  \lean{LionsC12.domainClauses_of_params, LionsC12.exists_eps, LionsC12.lions_counterexample_open,
    LionsC12.swf, LionsC12.dU'_eq, LionsC12.hasDerivAt_ofH10_dU, LionsC12.continuous_dU,
    LionsC12.continuous_dU', LionsC12.continuous_dF, LionsC12.not_holder_dB,
    LionsC12.isLionsSolutionΩ_dU, LionsC12.not_exists_L2_dU', LionsC12.not_exists_L2_form}
  \leanok
  \uses{def:domain-construction, lem:tensor, lem:weak-line, prop:coefficient, prop:holder-sharp,
    lem:solution-regularity, prop:forcing, lem:lions-uniqueness-domain}
  For every open $\Omega$, every admissible $\eps$, radius $r\leq1$ and centre $p\in\Omega$ with
  $\overline B(p,r(\pi+1))\subseteq\Omega$, the rescaled construction satisfies every clause of
  Corollary~\ref{cor:bounded-domains}.
\end{lemma}

\begin{proof}
  \leanok
  As described above.
\end{proof}


\chapter{Comparison with endpoint hypotheses in the literature}
Let $X=\Leb^\infty(0,\pi)$.  The paper compares the construction with four sufficient conditions
for maximal $\Leb^2$-regularity at the endpoint and deduces their failure from
Theorem~\ref{thm:main}: if one of them held, the cited theorem would give maximal regularity.
The cited theorems are not formalised.  The formalisation instead proves the failure of every
hypothesis directly from the construction, which does not depend on the literature.

\begin{lemma}[The operator norm of a form increment]
  \label{lem:opnorm}
  \lean{LionsC12.formOpDist, LionsC12.enorm_form_sub_le, LionsC12.formOpDist_le_eLpNorm,
    LionsC12.exists_form_sub_ge_of_lt_eLpNorm, LionsC12.formOpDist_eq_eLpNorm,
    LionsC12.formOpDist_a, LionsC12.aX, LionsC12.norm_aX_sub}
  \leanok
  \uses{def:lions-solution, def:coefficient}
  For a bounded measurable coefficient on an interval $(a,b)$ with continuous increments,
  $\norm{\cA(t)-\cA(s)}_{\cL(V,V')}=\norm{\fra(t)-\fra(s)}_{\Leb^\infty(a,b)}$.
\end{lemma}

\textbf{Lean form.}  The paper uses only the inequality $\leq$.  The operator norm is
\texttt{formOpDist}, the supremum of
$\abs{\fra(t)[v,w]-\fra(s)[v,w]}/(\norm{v}_V\norm{w}_V)$ over $v,w\neq0$, with the graph norm of
$V$; no operator $V\to V'$ is constructed.  The path $t\mapsto\fra(t)\in X$ is \texttt{aX}, with
values in Mathlib's $\Leb^\infty$.

\begin{proof}
  \leanok
  The upper bound is H\"older's inequality.  For the lower bound, test with
  $v=\psi((x-x_0)/h)$, a dilated smooth bump at a point where the increment is close to its
  supremum, for which $\norm{\partial_xv}^2/\norm{v}_V^2\to1$ as $h\to0$.
\end{proof}

\begin{corollary}[Failure of the known sufficient hypotheses at the endpoint]
  \label{cor:endpoint}
  \lean{LionsC12.endpoint_hypotheses_fail, LionsC12.eVariationOn_aX_eq_top,
    LionsC12.not_exists_dier_bound, LionsC12.not_exists_dini_form,
    LionsC12.not_exists_dini_coeff, LionsC12.dini_form_lintegral_eq_top,
    LionsC12.dini_coeff_lintegral_eq_top, LionsC12.gagliardoCoeff, LionsC12.gagliardoForm,
    LionsC12.gagliardoForm_eq, LionsC12.gagliardoCoeff_eq_top,
    LionsC12.exists_piece_gagliardo_eq_top, LionsC12.aeSquare,
    LionsC12.not_exists_auscherEgert_bound, LionsC12.not_exists_auscherEgert_bound',
    LionsC12.iSup_lintegral_formOpDist_sq_div_le, LionsC12.exists_partition_local_smallness,
    LionsC12.form_self_re_equiv}
  \leanok
  \uses{lem:opnorm, prop:coefficient, prop:holder-sharp}
  For the coefficient of Theorem~\ref{thm:main}:
  \begin{enumerate}
    \item the form does not have bounded variation in the sense of Dier, and
      $\fra\notin BV([0,1];X)$;
    \item there is no non-decreasing $\omega:[0,1]\to[0,\infty)$ with
      $\abs{\langle(\cA(t)-\cA(s))v,w\rangle}\leq\omega(\abs{t-s})\norm{v}_V\norm{w}_V$ for
      $s,t\in[0,1]$ and $\int_0^1\omega(r)r^{-3/2}\dd r<\infty$ (Haak--Ouhabaz); in particular no
      such modulus controls $\norm{\fra(t)-\fra(s)}_X$;
    \item neither $\cA:[0,1]\to\cL(V,V')$ nor $\fra:[0,1]\to X$ is piecewise $\Sob^{1/2}$ on any
      finite partition (Achache--Ouhabaz);
    \item there is no $M<\infty$ with
      $\sup_{I}\abs{I}^{-1}\iint_{I\times I}\abs{\fra(t,x)-\fra(s,x)}^2/\abs{t-s}^2\dd s\dd t\leq M$
      for a.e.\ $x\in(0,\pi)$ (Auscher--Egert, condition (8)).
  \end{enumerate}
  Moreover the two side conditions of Achache--Ouhabaz hold: the local smallness
  $\sup_{t\in I}\int_I\norm{\cA(t)-\cA(s)}^2_{\cL(V,V')}\abs{t-s}^{-1}\dd s\leq C\eps^2\abs{I}$, and the
  uniform equivalence $\re\fra(t)[v,v]\simeq\norm{v}_V^2$.
\end{corollary}

\textbf{Lean form.}  The statement is \texttt{endpoint\_hypotheses\_fail}, for every
$0<\eps\leq1$.  (i) is \texttt{eVariationOn} of the path \texttt{aX} on $[0,1]$, equal to
$\infty$, and the absence of a function $g$, monotone on $[0,1]$, with Dier's estimate.  (ii)
reads ``$\int_0^1\omega(r)r^{-3/2}\dd r<\infty$'' as finiteness of the lower Lebesgue integral of
$\omega(r)r^{-3/2}$ on $(0,1)$; the stronger versions
\texttt{dini\_form\_lintegral\_eq\_top}, \texttt{dini\_coeff\_lintegral\_eq\_top} say that the
integral is $\infty$ for every admissible monotone $\omega$.  (iii) says that every partition
$0=t_0<\dots<t_N=1$ has a piece $(t_k,t_{k+1})$ on which the Gagliardo seminorm
$\iint\norm{\fra(t)-\fra(s)}_X^2\abs{t-s}^{-2}$ (and the same for $\cA$, Lemma~\ref{lem:opnorm}) is
infinite; only the seminorm of the $\Sob^{1/2}$ norm is used, and the piece is the one containing
the accumulation time $5/16$.  (iv) has the supremum over open intervals $(p,q)\subseteq[0,1]$
inside the almost-everywhere quantifier, as in the paper; \texttt{not\_exists\_auscherEgert\_bound'}
excludes even the weaker hypothesis in which the exceptional null set may depend on the interval.  The local smallness constant is
$C=4C_\beta^2$.  The form equivalence is stated with $\norm{\partial_xv}_H$ in place of
$\norm{v}_V$ (the two are equivalent by Poincar\'e's inequality, which is not formalised for
$V$; the upper bound in the graph norm is \texttt{norm\_form\_a\_le}); the operator square root of
the Kato property is not formalised.  The theorems of Dier, Haak--Ouhabaz, Achache--Ouhabaz and
Auscher--Egert are not formalised.

\begin{proof}
  \leanok
  This is the argument Lean states, not the paper's, which deduces (i)--(iv) from the cited
  theorems and Theorem~\ref{thm:main}.  On the middle third $J_j$ of a block, the increment of
  $\fra(\cdot,0)$ over a quarter period of the first mode is at least $\frac{13}{15}\eps K_{j,1}^{-1}$;
  $3K_{j,1}$ such steps fit in $J_j$, so each block contributes at least $\eps$ to the variation,
  which proves (i), and Dier's estimate would bound it.  For (ii), on the pieces
  $(h_{j+1},h_j]$ of the scales $h_j=\pi/(2K_{j,1}^2)$, monotonicity gives
  $\omega(r)r^{-3/2}\geq\omega(h_{j+1})h_j^{-3/2}$, each piece contributes at least $c/(j+2)$, and the
  harmonic series diverges.  For (iii) and (iv), on rectangles $T\times S\subseteq J_j\times J_j$
  where the phase of the first mode lies in $[\pi/4,3\pi/4]$ and in $[5\pi/4,7\pi/4]$, the
  squared difference quotient of $\sin K^2(\cdot-\tau_j)$ integrates to at least
  $K^2\abs{J_j}/(36\pi)$, which, after the factor $\eps^2K^{-2}$, is $\geq c\eps^2\ell_j$ per block
  on a set of length $\ell_j$, hence unbounded after division by $\abs{I}$ and infinite in sum; (iv)
  integrates over $x$ by Tonelli.  Local smallness is
  Proposition~\ref{prop:coefficient} with Lemma~\ref{lem:opnorm}.
\end{proof}


\chapter{Lions' theorem}
The paper uses Lions' theorem for uniqueness: the forms are measurable, uniformly bounded and
uniformly coercive, so the variational problem with zero initial value has at most one solution
in $\Leb^2(0,1;V)\cap\Sob^1(0,1;V')$.  This chapter records the part of Lions' theory that the
formalisation proves.  It is not in the paper, which cites it.

\section{The abstract theory on real Gelfand triples}

\begin{definition}[Gelfand triples, form families and variational solutions]
  \label{def:gelfand}
  \lean{LionsC12.GelfandTriple, LionsC12.IsFormFamily, LionsC12.IsVariationalSolution}
  \leanok
  Real Hilbert spaces $V\hookrightarrow H$ with a dense continuous injection, a strongly
  measurable family of uniformly bounded, uniformly coercive bilinear forms on $V$, and the
  variational formulation of $u'+\cA(t)u=f$, $u(0)=u_0$, tested against $\C^1$ functions
  $\varphi:[0,T]\to V$ with $\varphi(T)=0$.
\end{definition}

\begin{theorem}[Lions' theorem]
  \label{thm:lions}
  \lean{LionsC12.lions_projection, LionsC12.lions_existence, LionsC12.lions_magenes,
    LionsC12.lions_uniqueness, LionsC12.lions_theorem, LionsC12.variational_continuous}
  \leanok
  \uses{def:gelfand}
  For $f\in\Leb^2(0,T;V')$ and $u_0\in H$ the variational problem has a unique solution
  $u\in\Leb^2(0,T;V)$, and $u$ has a representative in $\C([0,T];H)$ with $u(0)=u_0$.
\end{theorem}

\begin{proof}
  \leanok
  Existence by Lions' projection theorem (a Lax--Milgram theorem for a form continuous on a
  Hilbert space and coercive on a pre-Hilbert subspace).  Uniqueness by the Lions--Magenes
  energy identity $\norm{u(t)}_H^2=\norm{u(0)}_H^2+2\int_0^t\langle u',u\rangle$, proved with
  Steklov averages.
\end{proof}

\begin{theorem}[Lions' theorem under a Gårding inequality]
  \label{thm:lions-garding}
  \lean{LionsC12.lions_existence_garding}
  \leanok
  \uses{thm:lions}
  If the forms satisfy the Gårding inequality
  $\alpha\norm{v}_V^2\leq a(t)[v,v]+\lambda\norm{v}_H^2$, $\lambda\geq0$, instead of coercivity, the
  variational problem still has a solution $u\in\Leb^2(0,T;V)$, with
  $\norm{u}_{\Leb^2(0,T;V)}\leq\mathrm{e}^{\lambda T}(\norm{f}_{\Leb^2(0,T;V')}+\norm{u_0}_H)/\min\{\frac12,\alpha\}$.
\end{theorem}

\begin{proof}
  \leanok
  The exponential shift $u(t)=\mathrm{e}^{\lambda t}\tilde u(t)$ reduces to coercive forms
  $a(t)+\lambda\langle\cdot,\cdot\rangle_H$ and Theorem~\ref{thm:lions}.
\end{proof}

\section{Uniqueness in the paper's complex spaces}

\begin{lemma}[Uniqueness of the Lions solution on the line]
  \label{lem:lions-uniqueness}
  \lean{LionsC12.IsLionsSol, LionsC12.reDual, LionsC12.realTriple, LionsC12.IsLionsSol.ae_eq,
    LionsC12.isLionsSolution_unique, LionsC12.val_injective, LionsC12.denseRange_val}
  \leanok
  \uses{def:lions-solution, thm:lions}
  Let $U\subseteq\R$ be open and $\fra$ measurable, bounded and non-negative.  Two Lions solutions
  of $w'+\cA(t)w=f$, $w(0)=0$, in the sense of Definition~\ref{def:lions-solution}, agree
  almost everywhere on $(0,1)$.
\end{lemma}

\textbf{Lean form.}  The complex triple $V\hookrightarrow H\hookrightarrow V'$ is realified: $V$
and $H$ become real Hilbert spaces with the real part of the inner product, and an
anti-linear functional $\ell$ becomes the real functional $\re\ell$ (\texttt{reDual}, which is
injective).  The difference of two solutions solves the homogeneous problem in the realified
triple; the Lions--Magenes identity and $\re\fra(t)[w,w]\geq0$ give $\norm{w(t)}_H^2\leq0$.  Only
non-negativity of the form is used, so uniqueness is proved under a weaker hypothesis than the
paper's uniform coercivity.  The density of $V$ in $H$ is proved from the density of
$\C_c^\infty(U)$ in $\Leb^2(U)$ (Mathlib).

\begin{proof}
  \leanok
  As in the Lean form above.
\end{proof}

\begin{lemma}[Uniqueness of the Lions solution on a domain]
  \label{lem:lions-uniqueness-domain}
  \lean{LionsC12.isLionsSolutionΩ_unique, LionsC12.valΩ_injective, LionsC12.denseRange_valΩ}
  \leanok
  \uses{def:spaces-domain, lem:lions-uniqueness}
  The same statement for an open $\Omega\subseteq\R^d$ and the forms
  $\frb(t)[v,w]=\int_\Omega\frb\,\nabla v\cdot\overline{\nabla w}$ with $\frb$ measurable, bounded and
  non-negative.
\end{lemma}

\begin{proof}
  \leanok
  The proof of Lemma~\ref{lem:lions-uniqueness}, with the abstract statement
  \texttt{IsLionsSol.ae\_eq} applied to $V_\Omega\hookrightarrow H_\Omega$.
\end{proof}

